\documentclass[12pt, a4paper]{article}
\usepackage[utf8]{inputenc}
\usepackage[T1]{fontenc}
\usepackage[brazil]{babel}
\usepackage{geometry}
\usepackage{enumitem}
\usepackage{titlesec}
\usepackage{hyperref}
\geometry{margin=2.5cm}

\title{\textbf{Conceitos Necessários --- MS901/MT861}\\
\large Integração de Machine Learning e Análise Numérica para EDPs}
\author{Estudo pessoal}
\date{}

\begin{document}

\maketitle

% -------------------------------------------------------
\section*{Equações Diferenciais Parciais (EDPs)}

\begin{itemize}
    \item Classificação de EDPs: hiperbólicas, parabólicas, elípticas
    \item Problema de Cauchy: condições iniciais e de contorno
    \item EDPs lineares e não lineares
    \item Operador não linear $\mathcal{N}(u)$
    \item Equação de advecção linear: $u_t + a u_x = 0$
    \item Equação de Burgers invíscida: $u_t + \left(\frac{u^2}{2}\right)_x = 0$
    \item Equação de Burgers viscosa (com termo difusivo $\epsilon \Delta u$)
    \item Equação de Buckley-Leverett (modelo de dois fluidos em meios porosos)
    \item Leis de conservação hiperbólicas em 1D e 2D
    \item Problema de Riemann: dados iniciais descontínuos
    \item Ondas de choque (\textit{shock waves})
    \item Ondas de rarefação (\textit{rarefaction fan})
    \item Condição de Rankine-Hugoniot (velocidade de propagação do choque)
    \item Dados iniciais suaves vs.\ dados iniciais descontínuos
\end{itemize}

% -------------------------------------------------------
\section*{Soluções Fracas e Teoria de Entropia}

\begin{itemize}
    \item Definição de solução fraca (formulação variacional)
    \item Funções teste $\phi \in C^1_c(\mathbb{R}^2 \times \mathbb{R}_+)$
    \item Espaços funcionais: $L^\infty$, $L^1$, $C^1_c$
    \item Integração por partes em domínios espaço-temporais
    \item Não unicidade de soluções fracas
    \item Condição de entropia de Kru\v{z}kov: $\partial_t |u - c| + (\partial_x + \partial_y) Q[u;c] \leq 0$
    \item Par entropia-fluxo $(η, q)$: convexidade de $\eta$ e relação $Dq = D\eta \cdot Df(u)$
    \item Dissipação de entropia: $\partial_t \eta(u) + \partial_x q(u) + \partial_y q(u) \leq 0$
    \item Formulação variacional de Dafermos para o problema de Riemann
    \item Minimização da taxa de variação da entropia total $H = \frac{d}{dt}\int \eta(u)\,dxdy$
    \item Solução admissível (solução de entropia)
\end{itemize}

% -------------------------------------------------------
\section*{Análise Numérica de EDPs}

\begin{itemize}
    \item Consistência de um método numérico
    \item Estabilidade de um método numérico
    \item Convergência (Teorema de Lax-Richtmyer: consistência + estabilidade $\Rightarrow$ convergência)
    \item Condição CFL (Courant-Friedrichs-Lewy)
    \item Método de diferenças finitas (explícito e implícito)
    \item Método de Lax-Friedrichs
    \item Método de Lax-Wendroff
    \item Erro de truncamento e ordem de acurácia
    \item Análise de Fourier (von Neumann) para estabilidade
    \item Difusão numérica (\textit{numerical diffusion}) --- erros de ordem par
    \item Dispersão numérica (\textit{numerical dispersion}) --- erros de ordem ímpar
    \item Viscosidade artificial como mecanismo de estabilização
    \item Métodos \textit{upwind}
    \item Métodos de volumes finitos
    \item Discretização por Monte Carlo de funcionais integrais
\end{itemize}

% -------------------------------------------------------
\section*{Teoria da Equação Modificada}

\begin{itemize}
    \item Conceito de equação modificada (\textit{modified equation})
    \item Análise de Hirt: expansão em série de Taylor do esquema numérico
    \item Identificação do termo de erro dominante (difusivo ou dispersivo)
    \item Resíduo como perturbação: $\mathcal{L}(\hat{u}) = \epsilon$
    \item Diagnóstico físico do resíduo: difusão vs.\ dispersão
    \item Equação modificada de uma PINN (``Equação Modificada Neural'')
    \item Efeito da função de ativação e profundidade da rede na equação modificada
\end{itemize}

% -------------------------------------------------------
\section*{Physics-Informed Neural Networks (PINNs)}

\begin{itemize}
    \item Arquitetura de uma PINN: entrada $(x,t)$, saída $u_\theta(x,t)$
    \item Diferenciação automática (\textit{automatic differentiation})
    \item Função de perda residual da EDP: $\mathcal{L}_f(u) = \frac{1}{N_f}\sum |f(x_f^i, t_f^i)|^2$
    \item Função de perda de condição inicial: $\mathcal{L}_{u_0}$
    \item Função de perda de condição de contorno: $\mathcal{L}_{bc}$
    \item Perda total e problema de otimização: $\theta^* = \arg\min_\theta \max_\phi \mathcal{L}$
    \item Pontos de colocação (\textit{collocation points})
    \item Parâmetro de regularização $\lambda$ (ponderação das parcelas da loss)
    \item PINN clássica vs.\ PINN fraca (\textit{weakPINN})
    \item Formulação PDE-informed (informada pela estrutura da EDP)
    \item Formulação fraca da PINN: uso de funções teste e funcional de entropia
    \item Funcional de Kru\v{z}kov discretizado: $\mathcal{L}_0(u, \phi, c)$
    \item Adição de termo difusivo $\epsilon \Delta u$ para estabilização
    \item PINN híbrida (weakPINN + difusão)
    \item Métricas de erro: ELF e EEL (comparação com soluções de referência)
    \item Hiperparâmetros: número de camadas, neurônios por camada, épocas, $\epsilon$
\end{itemize}

% -------------------------------------------------------
\section*{Otimização e Aprendizado de Máquina}

\begin{itemize}
    \item Descida do gradiente (\textit{gradient descent})
    \item Otimizador SGD (\textit{Stochastic Gradient Descent})
    \item Otimizador Adam
    \item Análise de otimizadores como sistemas dinâmicos discretos
    \item Relação entre passo de tempo do otimizador e métodos de integração numérica
    \item \textit{Backpropagation}
    \item Funções de ativação: $\tanh$, ReLU, sigmoid, etc.
    \item \textit{Universal Approximation Theorem} (redes como aproximadores funcionais)
    \item Regularização e generalização
    \item \textit{Overfitting} e \textit{underfitting}
\end{itemize}

% -------------------------------------------------------
\section*{Tópicos Avançados e Transversais}

\begin{itemize}
    \item Maldição da dimensionalidade (\textit{curse of dimensionality})
    \item PINNs como método \textit{meshless} em alta dimensão
    \item Distinção entre erro numérico (discretização) e incerteza epistêmica (dados)
    \item Cálculo de alto desempenho (HPC) com GPUs
    \item Equações de leis de conservação em mecânica dos fluidos (advecção, Euler, Navier-Stokes)
    \item Modelos de escoamento em meios porosos (Buckley-Leverett, Darcy)
    \item Normas funcionais: $L^1$, $L^2$, $L^\infty$
    \item Espaços de Sobolev $H^1$, $W^{k,p}$
    \item Redação de relatório técnico científico em \LaTeX
    \item Interpretação física de resultados computacionais
\end{itemize}

\end{document}
